\section{Expected Outcome}

The expected outcome of this study is the development of a computational framework for analyzing strain-dependent electrical conduction 
in deformable materials. The framework combines continuum mechanics and electrical transport into a unified numerical model, 
enabling the systematic investigation of electromechanical coupling.

At the modeling level, the work provides a consistent formulation of conductivity as a strain-dependent quantity, 
together with a numerical implementation for solving the resulting coupled system. 
This allows for the analysis of how mechanical deformation influences electrical potential distribution, 
current density, and overall system sensitivity.

From a computational perspective, the study delivers a flexible pipeline that can accommodate different forms of the 
strain–conductivity relation, ranging from phenomenological descriptions to physically informed or data-driven models. 
This makes the framework adaptable to different materials and configurations.

Beyond the specific application, the results contribute to the understanding of nonlinear multiphysics systems 
involving coupling between deformation and transport phenomena. 
Potential applications include piezoresistive sensing, structural health monitoring, and functional materials design.

The developed framework also provides a foundation for future extensions toward multiscale modeling 
and more detailed transport descriptions.